\documentclass[12pt,a4paper]{article}

\usepackage[utf8]{inputenc}
\usepackage[T1]{fontenc}
\usepackage{enumitem}
\usepackage{geometry}
\geometry{margin=2.5cm}

\title{Questionário – Procura Sistemática de Literatura}
\author{Departamento de Física \\ Universidade Eduardo Mondlane}
\date{}

\begin{document}
\maketitle

\section*{Instruções}
Assinale a opção correta em cada questão.

\begin{enumerate}[label=\textbf{\arabic*.}]

\item O que caracteriza uma revisão sistemática?
\begin{enumerate}[label=\alph*)]
    \item Opinião do autor baseada em artigos selecionados
    \item Revisão narrativa sem critérios definidos
    \item Processo estruturado, objetivo e reproduzível
    \item Análise de um único estudo relevante
\end{enumerate}

\item Qual é o principal objetivo da pesquisa sistemática de literatura?
\begin{enumerate}[label=\alph*)]
    \item Reduzir o número de artigos publicados
    \item Confirmar apenas estudos que apoiem a hipótese
    \item Identificar e sintetizar evidências científicas de forma rigorosa
    \item Substituir estudos experimentais
\end{enumerate}

\item Qual das opções \textbf{NÃO} é uma característica de uma revisão sistemática?
\begin{enumerate}[label=\alph*)]
    \item Critérios de elegibilidade claros
    \item Estratégia de pesquisa documentada
    \item Seleção subjetiva dos estudos
    \item Síntese estruturada dos resultados
\end{enumerate}

\item O framework PICO/ST é utilizado para:
\begin{enumerate}[label=\alph*)]
    \item Avaliar a qualidade estatística dos dados
    \item Publicar artigos científicos
    \item Estruturar a questão de pesquisa
    \item Selecionar a revista científica
\end{enumerate}

\item No framework PICO/ST, a letra ``P'' representa:
\begin{enumerate}[label=\alph*)]
    \item Período de tempo
    \item Problema estatístico
    \item População
    \item Publicação
\end{enumerate}

\item Qual elemento do PICO/ST corresponde ao resultado de interesse?
\begin{enumerate}[label=\alph*)]
    \item P – População
    \item I – Intervenção
    \item C – Comparador
    \item O – Desfecho
\end{enumerate}

\item Qual operador booleano é utilizado para combinar sinônimos?
\begin{enumerate}[label=\alph*)]
    \item AND
    \item NOT
    \item OR
    \item WITH
\end{enumerate}

\item O operador booleano AND é usado para:
\begin{enumerate}[label=\alph*)]
    \item Excluir termos irrelevantes
    \item Combinar diferentes conceitos
    \item Incluir sinônimos
    \item Traduzir palavras automaticamente
\end{enumerate}

\item Qual é a função do operador NOT?
\begin{enumerate}[label=\alph*)]
    \item Combinar conceitos relacionados
    \item Incluir sinônimos
    \item Excluir termos indesejados
    \item Ampliar a pesquisa
\end{enumerate}

\item O truncamento (*) permite:
\begin{enumerate}[label=\alph*)]
    \item Excluir palavras específicas
    \item Incluir variações de uma mesma palavra
    \item Limitar a pesquisa a um único termo
    \item Traduzir termos automaticamente
\end{enumerate}

\item Qual das opções é uma base de dados científica?
\begin{enumerate}[label=\alph*)]
    \item Wikipedia
    \item Facebook
    \item Scopus
    \item WhatsApp
\end{enumerate}

\item Qual base de dados é mais utilizada na área biomédica?
\begin{enumerate}[label=\alph*)]
    \item Web of Science
    \item Scopus
    \item Google Scholar
    \item PubMed
\end{enumerate}

\item A aplicação de filtros (ano, idioma, tipo de estudo) serve para:
\begin{enumerate}[label=\alph*)]
    \item Eliminar artigos relevantes
    \item Tornar a pesquisa aleatória
    \item Refinar e tornar a busca mais precisa
    \item Aumentar o viés da pesquisa
\end{enumerate}

\item Qual é o primeiro passo na pesquisa sistemática de literatura?
\begin{enumerate}[label=\alph*)]
    \item Aplicar filtros
    \item Extrair resultados
    \item Definir a questão de pesquisa
    \item Analisar os artigos
\end{enumerate}

\item O principal objetivo do exercício prático é:
\begin{enumerate}[label=\alph*)]
    \item Publicar um artigo científico
    \item Escolher uma revista
    \item Aplicar os conceitos de pesquisa sistemática
    \item Aprender formatação em \LaTeX
\end{enumerate}

\end{enumerate}

\end{document}
\documentclass[12pt,a4paper]{article}

\usepackage[utf8]{inputenc}
\usepackage[T1]{fontenc}
\usepackage{enumitem}
\usepackage{geometry}
\geometry{margin=2.5cm}

\title{Questionário – Procura Sistemática de Literatura}
\author{Departamento de Física \\ Universidade Eduardo Mondlane}
\date{}

\begin{document}
\maketitle

\section*{Instruções}
Assinale a opção correta em cada questão.

\begin{enumerate}[label=\textbf{\arabic*.}]

\item O que caracteriza uma revisão sistemática?
\begin{enumerate}[label=\alph*)]
    \item Opinião do autor baseada em artigos selecionados
    \item Revisão narrativa sem critérios definidos
    \item Processo estruturado, objetivo e reproduzível
    \item Análise de um único estudo relevante
\end{enumerate}

\item Qual é o principal objetivo da pesquisa sistemática de literatura?
\begin{enumerate}[label=\alph*)]
    \item Reduzir o número de artigos publicados
    \item Confirmar apenas estudos que apoiem a hipótese
    \item Identificar e sintetizar evidências científicas de forma rigorosa
    \item Substituir estudos experimentais
\end{enumerate}

\item Qual das opções \textbf{NÃO} é uma característica de uma revisão sistemática?
\begin{enumerate}[label=\alph*)]
    \item Critérios de elegibilidade claros
    \item Estratégia de pesquisa documentada
    \item Seleção subjetiva dos estudos
    \item Síntese estruturada dos resultados
\end{enumerate}

\item O framework PICO/ST é utilizado para:
\begin{enumerate}[label=\alph*)]
    \item Avaliar a qualidade estatística dos dados
    \item Publicar artigos científicos
    \item Estruturar a questão de pesquisa
    \item Selecionar a revista científica
\end{enumerate}

\item No framework PICO/ST, a letra ``P'' representa:
\begin{enumerate}[label=\alph*)]
    \item Período de tempo
    \item Problema estatístico
    \item População
    \item Publicação
\end{enumerate}

\item Qual elemento do PICO/ST corresponde ao resultado de interesse?
\begin{enumerate}[label=\alph*)]
    \item P – População
    \item I – Intervenção
    \item C – Comparador
    \item O – Desfecho
\end{enumerate}

\item Qual operador booleano é utilizado para combinar sinônimos?
\begin{enumerate}[label=\alph*)]
    \item AND
    \item NOT
    \item OR
    \item WITH
\end{enumerate}

\item O operador booleano AND é usado para:
\begin{enumerate}[label=\alph*)]
    \item Excluir termos irrelevantes
    \item Combinar diferentes conceitos
    \item Incluir sinônimos
    \item Traduzir palavras automaticamente
\end{enumerate}

\item Qual é a função do operador NOT?
\begin{enumerate}[label=\alph*)]
    \item Combinar conceitos relacionados
    \item Incluir sinônimos
    \item Excluir termos indesejados
    \item Ampliar a pesquisa
\end{enumerate}

\item O truncamento (*) permite:
\begin{enumerate}[label=\alph*)]
    \item Excluir palavras específicas
    \item Incluir variações de uma mesma palavra
    \item Limitar a pesquisa a um único termo
    \item Traduzir termos automaticamente
\end{enumerate}

\item Qual das opções é uma base de dados científica?
\begin{enumerate}[label=\alph*)]
    \item Wikipedia
    \item Facebook
    \item Scopus
    \item WhatsApp
\end{enumerate}

\item Qual base de dados é mais utilizada na área biomédica?
\begin{enumerate}[label=\alph*)]
    \item Web of Science
    \item Scopus
    \item Google Scholar
    \item PubMed
\end{enumerate}

\item A aplicação de filtros (ano, idioma, tipo de estudo) serve para:
\begin{enumerate}[label=\alph*)]
    \item Eliminar artigos relevantes
    \item Tornar a pesquisa aleatória
    \item Refinar e tornar a busca mais precisa
    \item Aumentar o viés da pesquisa
\end{enumerate}

\item Qual é o primeiro passo na pesquisa sistemática de literatura?
\begin{enumerate}[label=\alph*)]
    \item Aplicar filtros
    \item Extrair resultados
    \item Definir a questão de pesquisa
    \item Analisar os artigos
\end{enumerate}

\item O principal objetivo do exercício prático é:
\begin{enumerate}[label=\alph*)]
    \item Publicar um artigo científico
    \item Escolher uma revista
    \item Aplicar os conceitos de pesquisa sistemática
    \item Aprender formatação em \LaTeX
\end{enumerate}

\end{enumerate}

\end{document}
